\chapter{Длинное название главы, в которой мы смотрим на примеры того, как будут верстаться изображения и списки} \label{chapt2}

\section{Одиночное изображение} \label{sect2_1}

\begin{figure} [htb] 
  \label{img:image1}
  \centering
  \includegraphics [scale=0.27] {tex}
  \caption{Подпись к картинке.} 
\end{figure}

%\newpage
%============================================================================================================================
\section{Длинное название параграфа, в котором мы узнаём как сделать две картинки с общим номером и названием} \label{sect2_2}

А это две картинки под общим номером и названием:
\begin{figure} [h!] 
   \centering\begin{tabular}{c}
     \includegraphics [scale=0.25] {tex} \\ 
      а) \\
      $\:$ \\   
     \includegraphics [scale=0.25] {tex} \\
      б) \\
   \end{tabular}
   \caption{Подпись к рисунку.} 
   \label{FBtsCA2}
\end{figure}

%\newpage
%============================================================================================================================
\section{Пример вёрстки списоков} \label{sect2_3}

Содержание по пунктам (например, выводы)  можно отобразить вот так:
\begin{enumerate}
 \item Первый вывод.
 \item Следующий вывод.
 \item Очередной вывод.
\end{enumerate}

\clearpage